%%%%%%%%%%%%%%%%%%%%%%%%%%%%%%%%%%%%%%%%%%%%%%%%%%%%%%%%%%%%
% 6.9 GRAPH THEORY
% The Composition of Advanced Mathematics
%%%%%%%%%%%%%%%%%%%%%%%%%%%%%%%%%%%%%%%%%%%%%%%%%%%%%%%%%%%%

\section{Graph Theory}
\label{sec:graph-theory}

\prerequisite{
Sets, relations, functions, basic counting, discrete structures,
recurrence relations, matrices, and elementary proof techniques.
}

\learningobjectives{
By the end of this section, the reader should be able to:
\begin{itemize}
    \item define graphs, vertices, edges, and adjacency;
    \item distinguish simple graphs, multigraphs, pseudographs,
          directed graphs, and weighted graphs;
    \item compute vertex degrees and apply the handshaking lemma;
    \item work with paths, walks, trails, circuits, and cycles;
    \item determine connectedness and connected components;
    \item identify complete, bipartite, multipartite, regular, and
          planar graphs;
    \item work with graph complements and graph isomorphisms;
    \item represent graphs using adjacency and incidence matrices;
    \item understand subgraphs and induced subgraphs;
    \item identify trees, forests, and spanning trees;
    \item apply basic tree properties;
    \item understand Euler trails and Euler circuits;
    \item understand Hamiltonian paths and cycles;
    \item work with graph colorings and chromatic number;
    \item understand matchings and vertex covers;
    \item state Hall's Marriage Theorem;
    \item recognize graph connectivity parameters;
    \item understand planar embeddings and Euler's planar formula;
    \item recognize graph algorithms such as BFS and DFS;
    \item connect graph theory with networks, optimization, probability,
          algebra, and combinatorics.
\end{itemize}
}

%%%%%%%%%%%%%%%%%%%%%%%%%%%%%%%%%%%%%%%%%%%%%%%%%%%%%%%%%%%%
\subsection{What Is a Graph?}
\label{subsec:what-is-a-graph}
%%%%%%%%%%%%%%%%%%%%%%%%%%%%%%%%%%%%%%%%%%%%%%%%%%%%%%%%%%%%

A graph is a mathematical structure used to represent pairwise
relationships among objects.

\begin{definitionbox}{Graph}
A graph is an ordered pair

\[
\boxed{
G=(V,E),
}
\]

where \(V\) is a set of vertices and \(E\) is a collection of edges
joining pairs of vertices.
\end{definitionbox}

The set \(V\) is called the \emph{vertex set}.

The set \(E\) is called the \emph{edge set}.

The symbols

\[
|V|
\]

and

\[
|E|
\]

denote the number of vertices and edges.

%%%%%%%%%%%%%%%%%%%%%%%%%%%%%%%%%%%%%%%%%%%%%%%%%%%%%%%%%%%%
\subsection{Vertices and Edges}
\label{subsec:vertices-and-edges}
%%%%%%%%%%%%%%%%%%%%%%%%%%%%%%%%%%%%%%%%%%%%%%%%%%%%%%%%%%%%

A vertex represents an object.

An edge represents a relationship between two objects.

For example, if

\[
V=\{A,B,C,D\},
\]

and

\[
E=
\{
\{A,B\},
\{A,C\},
\{B,D\}
\},
\]

then the graph contains four vertices and three edges.

Two vertices joined by an edge are called adjacent.

%%%%%%%%%%%%%%%%%%%%%%%%%%%%%%%%%%%%%%%%%%%%%%%%%%%%%%%%%%%%
\subsection{Simple Graphs}
\label{subsec:simple-graphs}
%%%%%%%%%%%%%%%%%%%%%%%%%%%%%%%%%%%%%%%%%%%%%%%%%%%%%%%%%%%%

\begin{definitionbox}{Simple Graph}
A simple graph is an undirected graph with:
\begin{itemize}
    \item no loops;
    \item no multiple edges between the same pair of vertices.
\end{itemize}
\end{definitionbox}

Thus every edge in a simple graph is an unordered pair

\[
\{u,v\}
\]

with

\[
u\neq v.
\]

Most of this section focuses on finite simple graphs unless otherwise
stated.

%%%%%%%%%%%%%%%%%%%%%%%%%%%%%%%%%%%%%%%%%%%%%%%%%%%%%%%%%%%%
\subsection{Multigraphs and Pseudographs}
\label{subsec:multigraphs-pseudographs}
%%%%%%%%%%%%%%%%%%%%%%%%%%%%%%%%%%%%%%%%%%%%%%%%%%%%%%%%%%%%

A multigraph may contain multiple edges joining the same pair of
vertices.

A pseudograph may additionally contain loops.

A loop is an edge joining a vertex to itself.

These generalized graph models are useful when pairwise relationships
may occur more than once.

Examples include:

\begin{itemize}
    \item transportation systems with several routes between cities;
    \item electrical networks with parallel connections;
    \item communication systems with multiple channels.
\end{itemize}

%%%%%%%%%%%%%%%%%%%%%%%%%%%%%%%%%%%%%%%%%%%%%%%%%%%%%%%%%%%%
\subsection{Directed Graphs}
\label{subsec:directed-graphs}
%%%%%%%%%%%%%%%%%%%%%%%%%%%%%%%%%%%%%%%%%%%%%%%%%%%%%%%%%%%%

A directed graph, or digraph, assigns an orientation to each edge.

A directed edge is an ordered pair

\[
(u,v).
\]

It is commonly written

\[
\boxed{
u\longrightarrow v.
}
\]

The direction matters.

Thus,

\[
(u,v)\neq(v,u)
\]

in general.

Directed graphs model asymmetric relationships such as:

\begin{itemize}
    \item web links;
    \item prerequisite structures;
    \item one-way roads;
    \item information flow;
    \item dominance relations.
\end{itemize}

%%%%%%%%%%%%%%%%%%%%%%%%%%%%%%%%%%%%%%%%%%%%%%%%%%%%%%%%%%%%
\subsection{Weighted Graphs}
\label{subsec:weighted-graphs}
%%%%%%%%%%%%%%%%%%%%%%%%%%%%%%%%%%%%%%%%%%%%%%%%%%%%%%%%%%%%

A weighted graph assigns a numerical value to each edge.

A weight function may be written

\[
\boxed{
w:E\to\mathbb{R}.
}
\]

The value \(w(e)\) may represent:

\begin{itemize}
    \item distance;
    \item travel time;
    \item cost;
    \item capacity;
    \item reliability;
    \item energy.
\end{itemize}

Weighted graphs are fundamental in optimization and network analysis.

%%%%%%%%%%%%%%%%%%%%%%%%%%%%%%%%%%%%%%%%%%%%%%%%%%%%%%%%%%%%
\subsection{Graph Drawings}
\label{subsec:graph-drawings}
%%%%%%%%%%%%%%%%%%%%%%%%%%%%%%%%%%%%%%%%%%%%%%%%%%%%%%%%%%%%

A graph may be drawn by representing vertices as points and edges as
curves joining them.

The drawing is not the graph itself.

Two very different drawings may represent the same abstract graph.

For example, crossing edges in a drawing do not automatically create a
new vertex.

\begin{importantbox}
A graph is determined by its vertex-edge relationships, not by the
specific geometric placement used in a drawing.
\end{importantbox}

%%%%%%%%%%%%%%%%%%%%%%%%%%%%%%%%%%%%%%%%%%%%%%%%%%%%%%%%%%%%
\subsection{Degree of a Vertex}
\label{subsec:degree-of-vertex}
%%%%%%%%%%%%%%%%%%%%%%%%%%%%%%%%%%%%%%%%%%%%%%%%%%%%%%%%%%%%

\begin{definitionbox}{Degree}
The degree of a vertex \(v\), written

\[
\boxed{
\deg(v),
}
\]

is the number of edges incident with \(v\).
\end{definitionbox}

In a simple graph, this is also the number of neighbors of \(v\).

A vertex of degree \(0\) is called isolated.

A vertex of degree \(1\) is often called a leaf or pendant vertex,
especially in a tree.

%%%%%%%%%%%%%%%%%%%%%%%%%%%%%%%%%%%%%%%%%%%%%%%%%%%%%%%%%%%%
\subsection{Neighborhood}
\label{subsec:graph-neighborhood}
%%%%%%%%%%%%%%%%%%%%%%%%%%%%%%%%%%%%%%%%%%%%%%%%%%%%%%%%%%%%

The open neighborhood of \(v\) is

\[
\boxed{
N(v)
=
\{u\in V:\{u,v\}\in E\}.
}
\]

The closed neighborhood is

\[
\boxed{
N[v]
=
N(v)\cup\{v\}.
}
\]

For a simple graph,

\[
\boxed{
\deg(v)=|N(v)|.
}
\]

Neighborhoods are central in local graph analysis.

%%%%%%%%%%%%%%%%%%%%%%%%%%%%%%%%%%%%%%%%%%%%%%%%%%%%%%%%%%%%
\subsection{Handshaking Lemma}
\label{subsec:handshaking-lemma}
%%%%%%%%%%%%%%%%%%%%%%%%%%%%%%%%%%%%%%%%%%%%%%%%%%%%%%%%%%%%

\begin{theorem}[Handshaking Lemma]
For every finite undirected graph,

\[
\boxed{
\sum_{v\in V}\deg(v)
=
2|E|.
}
\]
\end{theorem}

Every edge contributes exactly \(2\) to the total degree sum because it
has two endpoints.

%%%%%%%%%%%%%%%%%%%%%%%%%%%%%%%%%%%%%%%%%%%%%%%%%%%%%%%%%%%%
\subsection{Worked Example: Handshaking Lemma}
\label{subsec:worked-handshaking-lemma}
%%%%%%%%%%%%%%%%%%%%%%%%%%%%%%%%%%%%%%%%%%%%%%%%%%%%%%%%%%%%

\begin{workedexamplebox}{Finding the Number of Edges}

Suppose a graph has vertex degrees

\[
4,3,3,2,2,2.
\]

The total degree is

\[
4+3+3+2+2+2=16.
\]

By the handshaking lemma,

\[
2|E|=16.
\]

\begin{answerbox}
\[
\boxed{
|E|=8.
}
\]
\end{answerbox}

\end{workedexamplebox}

%%%%%%%%%%%%%%%%%%%%%%%%%%%%%%%%%%%%%%%%%%%%%%%%%%%%%%%%%%%%
\subsection{Odd-Degree Vertices}
\label{subsec:odd-degree-vertices}
%%%%%%%%%%%%%%%%%%%%%%%%%%%%%%%%%%%%%%%%%%%%%%%%%%%%%%%%%%%%

A direct consequence of the handshaking lemma is:

\begin{theorem}
Every finite undirected graph has an even number of vertices of odd
degree.
\end{theorem}

Since

\[
\sum_{v\in V}\deg(v)
\]

is even, the number of odd summands must be even.

%%%%%%%%%%%%%%%%%%%%%%%%%%%%%%%%%%%%%%%%%%%%%%%%%%%%%%%%%%%%
\subsection{Regular Graphs}
\label{subsec:regular-graphs}
%%%%%%%%%%%%%%%%%%%%%%%%%%%%%%%%%%%%%%%%%%%%%%%%%%%%%%%%%%%%

\begin{definitionbox}{Regular Graph}
A graph is \(r\)-regular if

\[
\boxed{
\deg(v)=r
}
\]

for every vertex \(v\).
\end{definitionbox}

If \(G\) is \(r\)-regular on \(n\) vertices, then

\[
nr=2|E|.
\]

Therefore,

\[
\boxed{
|E|=\frac{nr}{2}.
}
\]

In particular, \(nr\) must be even.

%%%%%%%%%%%%%%%%%%%%%%%%%%%%%%%%%%%%%%%%%%%%%%%%%%%%%%%%%%%%
\subsection{Complete Graphs}
\label{subsec:complete-graphs}
%%%%%%%%%%%%%%%%%%%%%%%%%%%%%%%%%%%%%%%%%%%%%%%%%%%%%%%%%%%%

The complete graph on \(n\) vertices is denoted

\[
\boxed{
K_n.
}
\]

Every pair of distinct vertices is adjacent.

Therefore,

\[
\deg(v)=n-1
\]

for every vertex.

The number of edges is

\[
\boxed{
|E(K_n)|
=
\binom{n}{2}
=
\frac{n(n-1)}{2}.
}
\]

%%%%%%%%%%%%%%%%%%%%%%%%%%%%%%%%%%%%%%%%%%%%%%%%%%%%%%%%%%%%
\subsection{Bipartite Graphs}
\label{subsec:bipartite-graphs}
%%%%%%%%%%%%%%%%%%%%%%%%%%%%%%%%%%%%%%%%%%%%%%%%%%%%%%%%%%%%

\begin{definitionbox}{Bipartite Graph}
A graph is bipartite if its vertex set can be partitioned into two
disjoint sets

\[
V=X\cup Y
\]

such that every edge has one endpoint in \(X\) and one endpoint in
\(Y\).
\end{definitionbox}

There are no edges joining two vertices within the same part.

%%%%%%%%%%%%%%%%%%%%%%%%%%%%%%%%%%%%%%%%%%%%%%%%%%%%%%%%%%%%
\subsection{Complete Bipartite Graphs}
\label{subsec:complete-bipartite-graphs}
%%%%%%%%%%%%%%%%%%%%%%%%%%%%%%%%%%%%%%%%%%%%%%%%%%%%%%%%%%%%

The complete bipartite graph with part sizes \(m\) and \(n\) is denoted

\[
\boxed{
K_{m,n}.
}
\]

Every vertex in the first part is joined to every vertex in the second.

Therefore,

\[
\boxed{
|E(K_{m,n})|=mn.
}
\]

%%%%%%%%%%%%%%%%%%%%%%%%%%%%%%%%%%%%%%%%%%%%%%%%%%%%%%%%%%%%
\subsection{Characterization of Bipartite Graphs}
\label{subsec:bipartite-characterization}
%%%%%%%%%%%%%%%%%%%%%%%%%%%%%%%%%%%%%%%%%%%%%%%%%%%%%%%%%%%%

\begin{theorem}
A graph is bipartite if and only if it contains no odd cycle.
\end{theorem}

Thus the existence of an odd cycle is exactly the obstruction to
bipartiteness.

This theorem connects local cycle structure with a global vertex
partition.

%%%%%%%%%%%%%%%%%%%%%%%%%%%%%%%%%%%%%%%%%%%%%%%%%%%%%%%%%%%%
\subsection{Multipartite Graphs}
\label{subsec:multipartite-graphs}
%%%%%%%%%%%%%%%%%%%%%%%%%%%%%%%%%%%%%%%%%%%%%%%%%%%%%%%%%%%%

A graph is \(r\)-partite if its vertex set can be divided into \(r\)
parts so that no edge joins two vertices in the same part.

A complete \(r\)-partite graph contains every possible edge between
different parts.

Tur\'an graphs introduced in
Section~\ref{sec:extremal-combinatorics}
are balanced complete multipartite graphs.

%%%%%%%%%%%%%%%%%%%%%%%%%%%%%%%%%%%%%%%%%%%%%%%%%%%%%%%%%%%%
\subsection{Subgraphs}
\label{subsec:subgraphs}
%%%%%%%%%%%%%%%%%%%%%%%%%%%%%%%%%%%%%%%%%%%%%%%%%%%%%%%%%%%%

\begin{definitionbox}{Subgraph}
A graph \(H\) is a subgraph of \(G\) if

\[
V(H)\subseteq V(G)
\]

and

\[
E(H)\subseteq E(G).
\]
\end{definitionbox}

A subgraph may omit vertices, edges, or both.

%%%%%%%%%%%%%%%%%%%%%%%%%%%%%%%%%%%%%%%%%%%%%%%%%%%%%%%%%%%%
\subsection{Induced Subgraphs}
\label{subsec:induced-subgraphs}
%%%%%%%%%%%%%%%%%%%%%%%%%%%%%%%%%%%%%%%%%%%%%%%%%%%%%%%%%%%%

For

\[
S\subseteq V(G),
\]

the induced subgraph

\[
\boxed{
G[S]
}
\]

contains vertex set \(S\) and every edge of \(G\) whose endpoints both
lie in \(S\).

Thus an induced subgraph is determined entirely by its chosen vertex
set.

%%%%%%%%%%%%%%%%%%%%%%%%%%%%%%%%%%%%%%%%%%%%%%%%%%%%%%%%%%%%
\subsection{Spanning Subgraphs}
\label{subsec:spanning-subgraphs}
%%%%%%%%%%%%%%%%%%%%%%%%%%%%%%%%%%%%%%%%%%%%%%%%%%%%%%%%%%%%

A spanning subgraph has the same vertex set as the original graph:

\[
V(H)=V(G).
\]

It may contain only some of the edges.

Spanning trees are especially important examples.

%%%%%%%%%%%%%%%%%%%%%%%%%%%%%%%%%%%%%%%%%%%%%%%%%%%%%%%%%%%%
\subsection{Graph Complement}
\label{subsec:graph-complement}
%%%%%%%%%%%%%%%%%%%%%%%%%%%%%%%%%%%%%%%%%%%%%%%%%%%%%%%%%%%%

For a simple graph \(G\), its complement is denoted

\[
\boxed{
\overline{G}.
}
\]

Two distinct vertices are adjacent in \(\overline{G}\) exactly when
they are not adjacent in \(G\).

Thus,

\[
\boxed{
E(G)\cap E(\overline{G})=\varnothing.
}
\]

Also,

\[
\boxed{
|E(G)|+|E(\overline{G})|
=
\binom{|V|}{2}.
}
\]

%%%%%%%%%%%%%%%%%%%%%%%%%%%%%%%%%%%%%%%%%%%%%%%%%%%%%%%%%%%%
\subsection{Graph Isomorphism}
\label{subsec:graph-isomorphism}
%%%%%%%%%%%%%%%%%%%%%%%%%%%%%%%%%%%%%%%%%%%%%%%%%%%%%%%%%%%%

\begin{definitionbox}{Graph Isomorphism}
Graphs \(G\) and \(H\) are isomorphic if there exists a bijection

\[
f:V(G)\to V(H)
\]

such that

\[
\boxed{
\{u,v\}\in E(G)
\iff
\{f(u),f(v)\}\in E(H).
}
\]
\end{definitionbox}

An isomorphism preserves adjacency.

Isomorphic graphs have the same abstract structure even if they are
drawn differently.

%%%%%%%%%%%%%%%%%%%%%%%%%%%%%%%%%%%%%%%%%%%%%%%%%%%%%%%%%%%%
\subsection{Graph Invariants}
\label{subsec:graph-invariants}
%%%%%%%%%%%%%%%%%%%%%%%%%%%%%%%%%%%%%%%%%%%%%%%%%%%%%%%%%%%%

A graph invariant is a property preserved under isomorphism.

Examples include:

\begin{itemize}
    \item number of vertices;
    \item number of edges;
    \item degree sequence;
    \item number of connected components;
    \item cycle lengths;
    \item chromatic number;
    \item adjacency spectrum.
\end{itemize}

Different values of an invariant prove non-isomorphism.

Matching invariants do not necessarily prove isomorphism.

%%%%%%%%%%%%%%%%%%%%%%%%%%%%%%%%%%%%%%%%%%%%%%%%%%%%%%%%%%%%
\subsection{Degree Sequences}
\label{subsec:degree-sequences}
%%%%%%%%%%%%%%%%%%%%%%%%%%%%%%%%%%%%%%%%%%%%%%%%%%%%%%%%%%%%

The degree sequence of a graph is the list of vertex degrees, typically
written in nonincreasing order.

For example,

\[
\boxed{
(4,3,3,2,2,0)
}
\]

is a possible degree sequence.

A natural question is whether a given integer sequence is
\emph{graphical}, meaning that some simple graph has exactly that
degree sequence.

%%%%%%%%%%%%%%%%%%%%%%%%%%%%%%%%%%%%%%%%%%%%%%%%%%%%%%%%%%%%
\subsection{Havel--Hakimi Principle}
\label{subsec:havel-hakimi}
%%%%%%%%%%%%%%%%%%%%%%%%%%%%%%%%%%%%%%%%%%%%%%%%%%%%%%%%%%%%

The Havel--Hakimi procedure tests graphical degree sequences.

Suppose

\[
d_1\geq d_2\geq\cdots\geq d_n.
\]

Remove \(d_1\) and subtract \(1\) from the next \(d_1\) entries.

The original sequence is graphical if and only if the resulting
sequence is graphical, provided the reduction remains valid.

This gives a recursive algorithmic test.

%%%%%%%%%%%%%%%%%%%%%%%%%%%%%%%%%%%%%%%%%%%%%%%%%%%%%%%%%%%%
\subsection{Walks}
\label{subsec:graph-walks}
%%%%%%%%%%%%%%%%%%%%%%%%%%%%%%%%%%%%%%%%%%%%%%%%%%%%%%%%%%%%

A walk is a sequence

\[
v_0,v_1,\ldots,v_k
\]

such that consecutive vertices are adjacent.

The length of the walk is

\[
k.
\]

Vertices and edges may repeat.

%%%%%%%%%%%%%%%%%%%%%%%%%%%%%%%%%%%%%%%%%%%%%%%%%%%%%%%%%%%%
\subsection{Trails}
\label{subsec:graph-trails}
%%%%%%%%%%%%%%%%%%%%%%%%%%%%%%%%%%%%%%%%%%%%%%%%%%%%%%%%%%%%

A trail is a walk in which no edge is repeated.

Vertices may still repeat.

Thus,

\[
\boxed{
\text{trail}
\Longrightarrow
\text{walk}.
}
\]

The converse need not hold.

%%%%%%%%%%%%%%%%%%%%%%%%%%%%%%%%%%%%%%%%%%%%%%%%%%%%%%%%%%%%
\subsection{Paths}
\label{subsec:graph-paths}
%%%%%%%%%%%%%%%%%%%%%%%%%%%%%%%%%%%%%%%%%%%%%%%%%%%%%%%%%%%%

A path is a walk in which no vertex is repeated.

A path with \(k\) edges has \(k+1\) vertices.

If a path begins at \(u\) and ends at \(v\), it is called a
\(u\)-\(v\) path.

%%%%%%%%%%%%%%%%%%%%%%%%%%%%%%%%%%%%%%%%%%%%%%%%%%%%%%%%%%%%
\subsection{Closed Walks and Circuits}
\label{subsec:closed-walks-circuits}
%%%%%%%%%%%%%%%%%%%%%%%%%%%%%%%%%%%%%%%%%%%%%%%%%%%%%%%%%%%%

A closed walk begins and ends at the same vertex.

A circuit is a closed trail.

A cycle is a closed path except that the first and last vertex are the
same.

Thus a cycle contains no repeated vertices other than its initial/final
vertex.

%%%%%%%%%%%%%%%%%%%%%%%%%%%%%%%%%%%%%%%%%%%%%%%%%%%%%%%%%%%%
\subsection{Cycles}
\label{subsec:graph-cycles}
%%%%%%%%%%%%%%%%%%%%%%%%%%%%%%%%%%%%%%%%%%%%%%%%%%%%%%%%%%%%

The cycle graph on \(n\) vertices is denoted

\[
\boxed{
C_n.
}
\]

Every vertex in \(C_n\) has degree \(2\).

The graph contains exactly one cycle using all \(n\) vertices.

%%%%%%%%%%%%%%%%%%%%%%%%%%%%%%%%%%%%%%%%%%%%%%%%%%%%%%%%%%%%
\subsection{Distance}
\label{subsec:graph-distance}
%%%%%%%%%%%%%%%%%%%%%%%%%%%%%%%%%%%%%%%%%%%%%%%%%%%%%%%%%%%%

The distance between vertices \(u\) and \(v\), written

\[
\boxed{
d(u,v),
}
\]

is the length of a shortest \(u\)-\(v\) path.

For vertices in the same connected component, graph distance defines a
metric.

Thus,

\[
d(u,v)\geq0,
\]

\[
d(u,v)=d(v,u),
\]

and

\[
d(u,w)\leq d(u,v)+d(v,w).
\]

%%%%%%%%%%%%%%%%%%%%%%%%%%%%%%%%%%%%%%%%%%%%%%%%%%%%%%%%%%%%
\subsection{Eccentricity}
\label{subsec:graph-eccentricity}
%%%%%%%%%%%%%%%%%%%%%%%%%%%%%%%%%%%%%%%%%%%%%%%%%%%%%%%%%%%%

For a connected graph, the eccentricity of \(v\) is

\[
\boxed{
\operatorname{ecc}(v)
=
\max_{u\in V}d(u,v).
}
\]

This measures the greatest distance from \(v\) to any other vertex.

%%%%%%%%%%%%%%%%%%%%%%%%%%%%%%%%%%%%%%%%%%%%%%%%%%%%%%%%%%%%
\subsection{Diameter and Radius}
\label{subsec:diameter-radius}
%%%%%%%%%%%%%%%%%%%%%%%%%%%%%%%%%%%%%%%%%%%%%%%%%%%%%%%%%%%%

The diameter of a connected graph is

\[
\boxed{
\operatorname{diam}(G)
=
\max_{u,v\in V}d(u,v).
}
\]

The radius is

\[
\boxed{
\operatorname{rad}(G)
=
\min_{v\in V}\operatorname{ecc}(v).
}
\]

Vertices attaining the radius form the center of the graph.

%%%%%%%%%%%%%%%%%%%%%%%%%%%%%%%%%%%%%%%%%%%%%%%%%%%%%%%%%%%%
\subsection{Connected Graphs}
\label{subsec:connected-graphs}
%%%%%%%%%%%%%%%%%%%%%%%%%%%%%%%%%%%%%%%%%%%%%%%%%%%%%%%%%%%%

\begin{definitionbox}{Connected Graph}
A graph is connected if every pair of vertices is joined by a path.
\end{definitionbox}

If no such path exists for some pair, the graph is disconnected.

Connectivity is one of the most fundamental structural properties of a
graph.

%%%%%%%%%%%%%%%%%%%%%%%%%%%%%%%%%%%%%%%%%%%%%%%%%%%%%%%%%%%%
\subsection{Connected Components}
\label{subsec:connected-components}
%%%%%%%%%%%%%%%%%%%%%%%%%%%%%%%%%%%%%%%%%%%%%%%%%%%%%%%%%%%%

A connected component is a maximal connected subgraph.

Every graph decomposes uniquely into connected components.

Thus,

\[
\boxed{
G
=
G_1\cup G_2\cup\cdots\cup G_k
}
\]

where each \(G_i\) is connected and the components have disjoint vertex
sets.

%%%%%%%%%%%%%%%%%%%%%%%%%%%%%%%%%%%%%%%%%%%%%%%%%%%%%%%%%%%%
\subsection{Cut Vertices}
\label{subsec:cut-vertices}
%%%%%%%%%%%%%%%%%%%%%%%%%%%%%%%%%%%%%%%%%%%%%%%%%%%%%%%%%%%%

A cut vertex is a vertex whose removal increases the number of connected
components.

Cut vertices represent structurally critical points in a network.

They are also called articulation vertices.

%%%%%%%%%%%%%%%%%%%%%%%%%%%%%%%%%%%%%%%%%%%%%%%%%%%%%%%%%%%%
\subsection{Bridges}
\label{subsec:bridges}
%%%%%%%%%%%%%%%%%%%%%%%%%%%%%%%%%%%%%%%%%%%%%%%%%%%%%%%%%%%%

A bridge is an edge whose removal increases the number of connected
components.

Equivalently, an edge is a bridge if and only if it lies on no cycle.

Thus cycles provide alternate routes around edges.

%%%%%%%%%%%%%%%%%%%%%%%%%%%%%%%%%%%%%%%%%%%%%%%%%%%%%%%%%%%%
\subsection{Vertex and Edge Connectivity}
\label{subsec:vertex-edge-connectivity}
%%%%%%%%%%%%%%%%%%%%%%%%%%%%%%%%%%%%%%%%%%%%%%%%%%%%%%%%%%%%

The vertex connectivity of a graph is denoted

\[
\boxed{
\kappa(G).
}
\]

It is the minimum number of vertices whose removal disconnects the graph
or reduces it to a trivial graph.

The edge connectivity is denoted

\[
\boxed{
\lambda(G).
}
\]

It is the minimum number of edges whose removal disconnects the graph.

For many nontrivial connected graphs,

\[
\boxed{
\kappa(G)\leq\lambda(G)\leq\delta(G),
}
\]

where

\[
\delta(G)
\]

is the minimum degree.

%%%%%%%%%%%%%%%%%%%%%%%%%%%%%%%%%%%%%%%%%%%%%%%%%%%%%%%%%%%%
\subsection{Adjacency Matrices}
\label{subsec:adjacency-matrices}
%%%%%%%%%%%%%%%%%%%%%%%%%%%%%%%%%%%%%%%%%%%%%%%%%%%%%%%%%%%%

Suppose

\[
V=\{v_1,\ldots,v_n\}.
\]

The adjacency matrix of a simple graph is

\[
A=(a_{ij}),
\]

where

\[
\boxed{
a_{ij}
=
\begin{cases}
1,&v_i\sim v_j,\\
0,&\text{otherwise}.
\end{cases}
}
\]

Here,

\[
v_i\sim v_j
\]

means \(v_i\) is adjacent to \(v_j\).

%%%%%%%%%%%%%%%%%%%%%%%%%%%%%%%%%%%%%%%%%%%%%%%%%%%%%%%%%%%%
\subsection{Properties of Adjacency Matrices}
\label{subsec:adjacency-matrix-properties}
%%%%%%%%%%%%%%%%%%%%%%%%%%%%%%%%%%%%%%%%%%%%%%%%%%%%%%%%%%%%

For a finite simple undirected graph:

\begin{itemize}
    \item \(A\) is symmetric;
    \item diagonal entries are \(0\);
    \item row sums equal vertex degrees.
\end{itemize}

Thus,

\[
\boxed{
\deg(v_i)
=
\sum_{j=1}^{n}a_{ij}.
}
\]

%%%%%%%%%%%%%%%%%%%%%%%%%%%%%%%%%%%%%%%%%%%%%%%%%%%%%%%%%%%%
\subsection{Powers of the Adjacency Matrix}
\label{subsec:powers-adjacency-matrix}
%%%%%%%%%%%%%%%%%%%%%%%%%%%%%%%%%%%%%%%%%%%%%%%%%%%%%%%%%%%%

A fundamental fact is:

\[
\boxed{
(A^k)_{ij}
=
\text{number of walks of length }k
\text{ from }v_i\text{ to }v_j.
}
\]

This creates a powerful connection between graph theory and linear
algebra.

%%%%%%%%%%%%%%%%%%%%%%%%%%%%%%%%%%%%%%%%%%%%%%%%%%%%%%%%%%%%
\subsection{Worked Example: Adjacency Matrix}
\label{subsec:worked-adjacency-matrix}
%%%%%%%%%%%%%%%%%%%%%%%%%%%%%%%%%%%%%%%%%%%%%%%%%%%%%%%%%%%%

\begin{workedexamplebox}{A Three-Vertex Path}

Consider

\[
P_3:
\qquad
v_1-v_2-v_3.
\]

Its adjacency matrix is

\[
A=
\begin{pmatrix}
0&1&0\\
1&0&1\\
0&1&0
\end{pmatrix}.
\]

Then

\[
A^2=
\begin{pmatrix}
1&0&1\\
0&2&0\\
1&0&1
\end{pmatrix}.
\]

The entry

\[
(A^2)_{13}=1
\]

shows there is exactly one walk of length \(2\) from \(v_1\) to \(v_3\).

\begin{answerbox}
\[
\boxed{
v_1\to v_2\to v_3
}
\]

is the unique length-\(2\) walk from \(v_1\) to \(v_3\).
\end{answerbox}

\end{workedexamplebox}

%%%%%%%%%%%%%%%%%%%%%%%%%%%%%%%%%%%%%%%%%%%%%%%%%%%%%%%%%%%%
\subsection{Incidence Matrices}
\label{subsec:graph-incidence-matrices}
%%%%%%%%%%%%%%%%%%%%%%%%%%%%%%%%%%%%%%%%%%%%%%%%%%%%%%%%%%%%

Suppose

\[
V=\{v_1,\ldots,v_n\}
\]

and

\[
E=\{e_1,\ldots,e_m\}.
\]

The incidence matrix

\[
B=(b_{ij})
\]

has

\[
\boxed{
b_{ij}
=
\begin{cases}
1,&v_i\text{ is incident with }e_j,\\
0,&\text{otherwise}.
\end{cases}
}
\]

For a simple undirected graph, each column contains exactly two \(1\)'s.

%%%%%%%%%%%%%%%%%%%%%%%%%%%%%%%%%%%%%%%%%%%%%%%%%%%%%%%%%%%%
\subsection{Trees}
\label{subsec:trees}
%%%%%%%%%%%%%%%%%%%%%%%%%%%%%%%%%%%%%%%%%%%%%%%%%%%%%%%%%%%%

\begin{definitionbox}{Tree}
A tree is a connected graph containing no cycles.
\end{definitionbox}

Trees are among the most important graph structures.

They describe hierarchy, branching, and minimal connectivity.

%%%%%%%%%%%%%%%%%%%%%%%%%%%%%%%%%%%%%%%%%%%%%%%%%%%%%%%%%%%%
\subsection{Equivalent Characterizations of Trees}
\label{subsec:tree-characterizations}
%%%%%%%%%%%%%%%%%%%%%%%%%%%%%%%%%%%%%%%%%%%%%%%%%%%%%%%%%%%%

For a finite graph \(T\) with \(n\) vertices, the following are
equivalent:

\begin{itemize}
    \item \(T\) is a tree;
    \item \(T\) is connected and acyclic;
    \item \(T\) is connected and has \(n-1\) edges;
    \item \(T\) is acyclic and has \(n-1\) edges;
    \item there is exactly one path between every pair of vertices;
    \item every edge is a bridge.
\end{itemize}

These equivalent descriptions make trees unusually flexible to study.

%%%%%%%%%%%%%%%%%%%%%%%%%%%%%%%%%%%%%%%%%%%%%%%%%%%%%%%%%%%%
\subsection{Number of Edges in a Tree}
\label{subsec:tree-edge-count}
%%%%%%%%%%%%%%%%%%%%%%%%%%%%%%%%%%%%%%%%%%%%%%%%%%%%%%%%%%%%

\begin{theorem}
Every tree with \(n\) vertices has exactly

\[
\boxed{
n-1
}
\]

edges.
\end{theorem}

This may be proved by induction using the existence of leaves.

%%%%%%%%%%%%%%%%%%%%%%%%%%%%%%%%%%%%%%%%%%%%%%%%%%%%%%%%%%%%
\subsection{Leaves in Trees}
\label{subsec:tree-leaves}
%%%%%%%%%%%%%%%%%%%%%%%%%%%%%%%%%%%%%%%%%%%%%%%%%%%%%%%%%%%%

Every finite tree with at least two vertices has at least two leaves.

Removing a leaf and its incident edge from a tree produces another
tree.

This property makes induction particularly natural for tree proofs.

%%%%%%%%%%%%%%%%%%%%%%%%%%%%%%%%%%%%%%%%%%%%%%%%%%%%%%%%%%%%
\subsection{Forests}
\label{subsec:forests}
%%%%%%%%%%%%%%%%%%%%%%%%%%%%%%%%%%%%%%%%%%%%%%%%%%%%%%%%%%%%

A forest is an acyclic graph.

Each connected component of a forest is a tree.

If a forest has \(n\) vertices and \(c\) connected components, then

\[
\boxed{
|E|
=
n-c.
}
\]

%%%%%%%%%%%%%%%%%%%%%%%%%%%%%%%%%%%%%%%%%%%%%%%%%%%%%%%%%%%%
\subsection{Spanning Trees}
\label{subsec:spanning-trees}
%%%%%%%%%%%%%%%%%%%%%%%%%%%%%%%%%%%%%%%%%%%%%%%%%%%%%%%%%%%%

\begin{definitionbox}{Spanning Tree}
A spanning tree of a connected graph \(G\) is a subgraph \(T\) such that:
\begin{itemize}
    \item \(V(T)=V(G)\);
    \item \(T\) is a tree.
\end{itemize}
\end{definitionbox}

A spanning tree preserves connectivity while using the minimum possible
number of edges.

%%%%%%%%%%%%%%%%%%%%%%%%%%%%%%%%%%%%%%%%%%%%%%%%%%%%%%%%%%%%
\subsection{Every Connected Graph Has a Spanning Tree}
\label{subsec:connected-spanning-tree}
%%%%%%%%%%%%%%%%%%%%%%%%%%%%%%%%%%%%%%%%%%%%%%%%%%%%%%%%%%%%

Take a connected graph.

If it contains a cycle, remove one edge from that cycle.

Connectivity is preserved.

Repeat until no cycles remain.

The resulting spanning subgraph is connected and acyclic.

Therefore it is a spanning tree.

%%%%%%%%%%%%%%%%%%%%%%%%%%%%%%%%%%%%%%%%%%%%%%%%%%%%%%%%%%%%
\subsection{Cayley's Formula}
\label{subsec:cayleys-formula}
%%%%%%%%%%%%%%%%%%%%%%%%%%%%%%%%%%%%%%%%%%%%%%%%%%%%%%%%%%%%

\begin{theorem}[Cayley's Formula]
The number of labeled trees on \(n\) vertices is

\[
\boxed{
n^{n-2}.
}
\]
\end{theorem}

For example, the number of labeled trees on \(4\) vertices is

\[
4^{4-2}=16.
\]

One classical proof uses Prüfer codes.

%%%%%%%%%%%%%%%%%%%%%%%%%%%%%%%%%%%%%%%%%%%%%%%%%%%%%%%%%%%%
\subsection{Prüfer Codes}
\label{subsec:prufer-codes}
%%%%%%%%%%%%%%%%%%%%%%%%%%%%%%%%%%%%%%%%%%%%%%%%%%%%%%%%%%%%

A labeled tree on vertex set

\[
[n]
\]

can be encoded by a sequence of length

\[
n-2.
\]

Each sequence entry belongs to \([n]\).

Thus there are

\[
n^{n-2}
\]

possible Prüfer sequences.

The Prüfer correspondence is bijective, giving Cayley's formula.

%%%%%%%%%%%%%%%%%%%%%%%%%%%%%%%%%%%%%%%%%%%%%%%%%%%%%%%%%%%%
\subsection{Rooted Trees}
\label{subsec:rooted-trees}
%%%%%%%%%%%%%%%%%%%%%%%%%%%%%%%%%%%%%%%%%%%%%%%%%%%%%%%%%%%%

A rooted tree has one distinguished vertex called the root.

This creates hierarchical language such as:

\begin{itemize}
    \item parent;
    \item child;
    \item ancestor;
    \item descendant;
    \item depth;
    \item subtree.
\end{itemize}

Rooted trees are fundamental in computer science and recursive
combinatorics.

%%%%%%%%%%%%%%%%%%%%%%%%%%%%%%%%%%%%%%%%%%%%%%%%%%%%%%%%%%%%
\subsection{Euler Trails}
\label{subsec:euler-trails}
%%%%%%%%%%%%%%%%%%%%%%%%%%%%%%%%%%%%%%%%%%%%%%%%%%%%%%%%%%%%

An Euler trail uses every edge exactly once.

An Euler circuit is an Euler trail that begins and ends at the same
vertex.

These questions were historically central to the origins of graph
theory.

%%%%%%%%%%%%%%%%%%%%%%%%%%%%%%%%%%%%%%%%%%%%%%%%%%%%%%%%%%%%
\subsection{Euler Circuit Criterion}
\label{subsec:euler-circuit-criterion}
%%%%%%%%%%%%%%%%%%%%%%%%%%%%%%%%%%%%%%%%%%%%%%%%%%%%%%%%%%%%

\begin{theorem}[Euler Circuit Criterion]
A connected graph has an Euler circuit if and only if every vertex has
even degree.
\end{theorem}

Thus,

\[
\boxed{
\text{Euler circuit}
\iff
\text{all degrees even}.
}
\]

%%%%%%%%%%%%%%%%%%%%%%%%%%%%%%%%%%%%%%%%%%%%%%%%%%%%%%%%%%%%
\subsection{Euler Trail Criterion}
\label{subsec:euler-trail-criterion}
%%%%%%%%%%%%%%%%%%%%%%%%%%%%%%%%%%%%%%%%%%%%%%%%%%%%%%%%%%%%

\begin{theorem}[Euler Trail Criterion]
A connected graph has an Euler trail but no Euler circuit if and only if
exactly two vertices have odd degree.
\end{theorem}

The trail begins at one odd-degree vertex and ends at the other.

%%%%%%%%%%%%%%%%%%%%%%%%%%%%%%%%%%%%%%%%%%%%%%%%%%%%%%%%%%%%
\subsection{Worked Example: Euler Trail}
\label{subsec:worked-euler-trail}
%%%%%%%%%%%%%%%%%%%%%%%%%%%%%%%%%%%%%%%%%%%%%%%%%%%%%%%%%%%%

\begin{workedexamplebox}{Checking Euler Structure}

Suppose a connected graph has degrees

\[
4,4,2,2,3,3.
\]

Exactly two vertices have odd degree.

Therefore the graph has an Euler trail.

Because not every degree is even, it does not have an Euler circuit.

\begin{answerbox}
\[
\boxed{
\text{Euler trail: yes}
\qquad
\text{Euler circuit: no}.
}
\]
\end{answerbox}

\end{workedexamplebox}

%%%%%%%%%%%%%%%%%%%%%%%%%%%%%%%%%%%%%%%%%%%%%%%%%%%%%%%%%%%%
\subsection{Hamiltonian Paths and Cycles}
\label{subsec:hamiltonian-paths-cycles}
%%%%%%%%%%%%%%%%%%%%%%%%%%%%%%%%%%%%%%%%%%%%%%%%%%%%%%%%%%%%

A Hamiltonian path visits every vertex exactly once.

A Hamiltonian cycle visits every vertex exactly once and returns to the
starting vertex.

Unlike Euler problems, Hamiltonian problems concern vertices rather than
edges.

There is no simple degree characterization comparable to Euler's
criterion.

%%%%%%%%%%%%%%%%%%%%%%%%%%%%%%%%%%%%%%%%%%%%%%%%%%%%%%%%%%%%
\subsection{Dirac's Theorem}
\label{subsec:diracs-theorem}
%%%%%%%%%%%%%%%%%%%%%%%%%%%%%%%%%%%%%%%%%%%%%%%%%%%%%%%%%%%%

\begin{theorem}[Dirac's Theorem]
Let \(G\) be a simple graph on \(n\geq3\) vertices.

If

\[
\boxed{
\deg(v)\geq\frac n2
}
\]

for every vertex \(v\), then \(G\) contains a Hamiltonian cycle.
\end{theorem}

This condition is sufficient but not necessary.

%%%%%%%%%%%%%%%%%%%%%%%%%%%%%%%%%%%%%%%%%%%%%%%%%%%%%%%%%%%%
\subsection{Ore's Theorem}
\label{subsec:ores-theorem}
%%%%%%%%%%%%%%%%%%%%%%%%%%%%%%%%%%%%%%%%%%%%%%%%%%%%%%%%%%%%

\begin{theorem}[Ore's Theorem]
Let \(G\) be a simple graph on \(n\geq3\) vertices.

If every pair of nonadjacent vertices \(u,v\) satisfies

\[
\boxed{
\deg(u)+\deg(v)\geq n,
}
\]

then \(G\) is Hamiltonian.
\end{theorem}

Dirac's theorem follows as a special case.

%%%%%%%%%%%%%%%%%%%%%%%%%%%%%%%%%%%%%%%%%%%%%%%%%%%%%%%%%%%%
\subsection{Graph Coloring}
\label{subsec:graph-coloring}
%%%%%%%%%%%%%%%%%%%%%%%%%%%%%%%%%%%%%%%%%%%%%%%%%%%%%%%%%%%%

A proper vertex coloring assigns colors to vertices so that adjacent
vertices receive different colors.

A \(k\)-coloring uses at most \(k\) colors.

%%%%%%%%%%%%%%%%%%%%%%%%%%%%%%%%%%%%%%%%%%%%%%%%%%%%%%%%%%%%
\subsection{Chromatic Number}
\label{subsec:chromatic-number}
%%%%%%%%%%%%%%%%%%%%%%%%%%%%%%%%%%%%%%%%%%%%%%%%%%%%%%%%%%%%

\begin{definitionbox}{Chromatic Number}
The chromatic number of a graph \(G\), denoted

\[
\boxed{
\chi(G),
}
\]

is the minimum number of colors required for a proper vertex coloring.
\end{definitionbox}

Examples include

\[
\chi(K_n)=n
\]

and

\[
\chi(C_{2m})=2.
\]

For odd cycles,

\[
\boxed{
\chi(C_{2m+1})=3.
}
\]

%%%%%%%%%%%%%%%%%%%%%%%%%%%%%%%%%%%%%%%%%%%%%%%%%%%%%%%%%%%%
\subsection{Bipartite Graphs and Coloring}
\label{subsec:bipartite-coloring}
%%%%%%%%%%%%%%%%%%%%%%%%%%%%%%%%%%%%%%%%%%%%%%%%%%%%%%%%%%%%

A nonempty graph is bipartite if and only if

\[
\boxed{
\chi(G)\leq2.
}
\]

Thus bipartiteness may be interpreted either structurally through a
vertex partition or algorithmically through two-colorability.

%%%%%%%%%%%%%%%%%%%%%%%%%%%%%%%%%%%%%%%%%%%%%%%%%%%%%%%%%%%%
\subsection{Greedy Coloring}
\label{subsec:greedy-coloring}
%%%%%%%%%%%%%%%%%%%%%%%%%%%%%%%%%%%%%%%%%%%%%%%%%%%%%%%%%%%%

Choose an ordering of the vertices.

Color them one at a time, always assigning the lowest available color.

If the maximum degree is

\[
\Delta(G),
\]

then each vertex has at most \(\Delta(G)\) colored neighbors.

Therefore,

\[
\boxed{
\chi(G)\leq\Delta(G)+1.
}
\]

This is the basic greedy coloring bound.

%%%%%%%%%%%%%%%%%%%%%%%%%%%%%%%%%%%%%%%%%%%%%%%%%%%%%%%%%%%%
\subsection{Chromatic Polynomial}
\label{subsec:chromatic-polynomial}
%%%%%%%%%%%%%%%%%%%%%%%%%%%%%%%%%%%%%%%%%%%%%%%%%%%%%%%%%%%%

The chromatic polynomial

\[
\boxed{
P_G(k)
}
\]

counts proper colorings of \(G\) using \(k\) available colors.

For a tree \(T\) with \(n\) vertices,

\[
\boxed{
P_T(k)
=
k(k-1)^{n-1}.
}
\]

For the complete graph,

\[
\boxed{
P_{K_n}(k)
=
k(k-1)\cdots(k-n+1).
}
\]

%%%%%%%%%%%%%%%%%%%%%%%%%%%%%%%%%%%%%%%%%%%%%%%%%%%%%%%%%%%%
\subsection{Deletion--Contraction}
\label{subsec:graph-deletion-contraction}
%%%%%%%%%%%%%%%%%%%%%%%%%%%%%%%%%%%%%%%%%%%%%%%%%%%%%%%%%%%%

For an edge \(e\),

\[
\boxed{
P_G(k)
=
P_{G-e}(k)
-
P_{G/e}(k).
}
\]

Here \(G-e\) deletes \(e\), while \(G/e\) contracts \(e\).

This recurrence connects graph coloring with algebraic combinatorics and
matroid theory.

%%%%%%%%%%%%%%%%%%%%%%%%%%%%%%%%%%%%%%%%%%%%%%%%%%%%%%%%%%%%
\subsection{Edge Coloring}
\label{subsec:edge-coloring}
%%%%%%%%%%%%%%%%%%%%%%%%%%%%%%%%%%%%%%%%%%%%%%%%%%%%%%%%%%%%

An edge coloring assigns colors to edges so that adjacent edges receive
different colors.

The edge chromatic number is denoted

\[
\boxed{
\chi'(G).
}
\]

It is the minimum number of colors needed for a proper edge coloring.

%%%%%%%%%%%%%%%%%%%%%%%%%%%%%%%%%%%%%%%%%%%%%%%%%%%%%%%%%%%%
\subsection{Vizing's Theorem}
\label{subsec:vizing-theorem}
%%%%%%%%%%%%%%%%%%%%%%%%%%%%%%%%%%%%%%%%%%%%%%%%%%%%%%%%%%%%

\begin{theorem}[Vizing's Theorem]
For every finite simple graph,

\[
\boxed{
\Delta(G)
\leq
\chi'(G)
\leq
\Delta(G)+1.
}
\]
\end{theorem}

Thus every simple graph belongs to one of two edge-coloring classes.

%%%%%%%%%%%%%%%%%%%%%%%%%%%%%%%%%%%%%%%%%%%%%%%%%%%%%%%%%%%%
\subsection{Matchings}
\label{subsec:matchings}
%%%%%%%%%%%%%%%%%%%%%%%%%%%%%%%%%%%%%%%%%%%%%%%%%%%%%%%%%%%%

\begin{definitionbox}{Matching}
A matching is a set of edges no two of which share a vertex.
\end{definitionbox}

A maximum matching is a matching having the largest possible number of
edges.

The matching number is often denoted

\[
\boxed{
\nu(G).
}
\]

%%%%%%%%%%%%%%%%%%%%%%%%%%%%%%%%%%%%%%%%%%%%%%%%%%%%%%%%%%%%
\subsection{Perfect Matchings}
\label{subsec:perfect-matchings}
%%%%%%%%%%%%%%%%%%%%%%%%%%%%%%%%%%%%%%%%%%%%%%%%%%%%%%%%%%%%

A perfect matching covers every vertex exactly once.

Thus a graph with a perfect matching must have an even number of
vertices.

Perfect matchings model complete pairings and assignments.

%%%%%%%%%%%%%%%%%%%%%%%%%%%%%%%%%%%%%%%%%%%%%%%%%%%%%%%%%%%%
\subsection{Bipartite Matching}
\label{subsec:bipartite-matching}
%%%%%%%%%%%%%%%%%%%%%%%%%%%%%%%%%%%%%%%%%%%%%%%%%%%%%%%%%%%%

Suppose

\[
G=(X\cup Y,E)
\]

is bipartite.

A matching pairs vertices from \(X\) with vertices from \(Y\).

Such problems model:

\begin{itemize}
    \item workers to jobs;
    \item students to projects;
    \item machines to tasks;
    \item flights to gates.
\end{itemize}

%%%%%%%%%%%%%%%%%%%%%%%%%%%%%%%%%%%%%%%%%%%%%%%%%%%%%%%%%%%%
\subsection{Hall's Marriage Theorem}
\label{subsec:halls-marriage-theorem}
%%%%%%%%%%%%%%%%%%%%%%%%%%%%%%%%%%%%%%%%%%%%%%%%%%%%%%%%%%%%

\begin{theorem}[Hall's Marriage Theorem]
Let

\[
G=(X\cup Y,E)
\]

be bipartite.

There exists a matching saturating every vertex of \(X\) if and only if

\[
\boxed{
|N(S)|\geq|S|
}
\]

for every subset

\[
S\subseteq X.
\]
\end{theorem}

Here,

\[
N(S)
\]

is the set of neighbors in \(Y\) adjacent to at least one vertex of
\(S\).

%%%%%%%%%%%%%%%%%%%%%%%%%%%%%%%%%%%%%%%%%%%%%%%%%%%%%%%%%%%%
\subsection{Worked Example: Hall's Condition}
\label{subsec:worked-hall}
%%%%%%%%%%%%%%%%%%%%%%%%%%%%%%%%%%%%%%%%%%%%%%%%%%%%%%%%%%%%

\begin{workedexamplebox}{Checking a Matching Condition}

Suppose a bipartite graph has left part

\[
X=\{a,b,c\}
\]

and right part

\[
Y=\{1,2,3,4\}.
\]

Assume

\[
N(\{a\})=\{1,2\},
\]

\[
N(\{b\})=\{2,3\},
\]

\[
N(\{c\})=\{3,4\},
\]

and every larger subset has at least as many neighbors as vertices.

Then Hall's condition holds.

\begin{answerbox}
\[
\boxed{
\text{There exists a matching saturating }X.
}
\]
\end{answerbox}

\end{workedexamplebox}

%%%%%%%%%%%%%%%%%%%%%%%%%%%%%%%%%%%%%%%%%%%%%%%%%%%%%%%%%%%%
\subsection{Vertex Covers}
\label{subsec:vertex-covers}
%%%%%%%%%%%%%%%%%%%%%%%%%%%%%%%%%%%%%%%%%%%%%%%%%%%%%%%%%%%%

\begin{definitionbox}{Vertex Cover}
A vertex cover is a set

\[
C\subseteq V
\]

such that every edge has at least one endpoint in \(C\).
\end{definitionbox}

The minimum size of a vertex cover is denoted

\[
\boxed{
\tau(G).
}
\]

%%%%%%%%%%%%%%%%%%%%%%%%%%%%%%%%%%%%%%%%%%%%%%%%%%%%%%%%%%%%
\subsection{Independent Sets}
\label{subsec:independent-sets}
%%%%%%%%%%%%%%%%%%%%%%%%%%%%%%%%%%%%%%%%%%%%%%%%%%%%%%%%%%%%

An independent set is a set of vertices containing no adjacent pair.

The maximum size of an independent set is denoted

\[
\boxed{
\alpha(G).
}
\]

A set \(C\) is a vertex cover if and only if

\[
V\setminus C
\]

is independent.

Therefore,

\[
\boxed{
\alpha(G)+\tau(G)=|V|.
}
\]

%%%%%%%%%%%%%%%%%%%%%%%%%%%%%%%%%%%%%%%%%%%%%%%%%%%%%%%%%%%%
\subsection{K\H{o}nig's Theorem}
\label{subsec:konigs-theorem}
%%%%%%%%%%%%%%%%%%%%%%%%%%%%%%%%%%%%%%%%%%%%%%%%%%%%%%%%%%%%

\begin{theorem}[K\H{o}nig's Theorem]
For every bipartite graph,

\[
\boxed{
\nu(G)=\tau(G).
}
\]
\end{theorem}

Thus the maximum size of a matching equals the minimum size of a vertex
cover.

This is a fundamental min-max theorem.

%%%%%%%%%%%%%%%%%%%%%%%%%%%%%%%%%%%%%%%%%%%%%%%%%%%%%%%%%%%%
\subsection{Planar Graphs}
\label{subsec:planar-graphs}
%%%%%%%%%%%%%%%%%%%%%%%%%%%%%%%%%%%%%%%%%%%%%%%%%%%%%%%%%%%%

A graph is planar if it can be drawn in the plane with no edge
crossings except at common endpoints.

A particular crossing-free drawing is called a planar embedding.

%%%%%%%%%%%%%%%%%%%%%%%%%%%%%%%%%%%%%%%%%%%%%%%%%%%%%%%%%%%%
\subsection{Faces of a Planar Graph}
\label{subsec:planar-faces}
%%%%%%%%%%%%%%%%%%%%%%%%%%%%%%%%%%%%%%%%%%%%%%%%%%%%%%%%%%%%

A planar embedding divides the plane into regions called faces.

One face is unbounded and is called the outer face.

If

\[
f
\]

denotes the number of faces, Euler's formula relates \(f\) to the
numbers of vertices and edges.

%%%%%%%%%%%%%%%%%%%%%%%%%%%%%%%%%%%%%%%%%%%%%%%%%%%%%%%%%%%%
\subsection{Euler's Planar Formula}
\label{subsec:eulers-planar-formula}
%%%%%%%%%%%%%%%%%%%%%%%%%%%%%%%%%%%%%%%%%%%%%%%%%%%%%%%%%%%%

\begin{theorem}[Euler's Formula]
For a connected planar graph,

\[
\boxed{
|V|-|E|+|F|=2.
}
\]
\end{theorem}

Here,

\[
|F|
\]

denotes the number of faces in a planar embedding.

%%%%%%%%%%%%%%%%%%%%%%%%%%%%%%%%%%%%%%%%%%%%%%%%%%%%%%%%%%%%
\subsection{Worked Example: Euler's Formula}
\label{subsec:worked-euler-planar}
%%%%%%%%%%%%%%%%%%%%%%%%%%%%%%%%%%%%%%%%%%%%%%%%%%%%%%%%%%%%

\begin{workedexamplebox}{Finding the Number of Faces}

Suppose a connected planar graph has

\[
|V|=8
\]

and

\[
|E|=12.
\]

Euler's formula gives

\[
8-12+|F|=2.
\]

Therefore,

\[
|F|=6.
\]

\begin{answerbox}
\[
\boxed{
|F|=6.
}
\]
\end{answerbox}

\end{workedexamplebox}

%%%%%%%%%%%%%%%%%%%%%%%%%%%%%%%%%%%%%%%%%%%%%%%%%%%%%%%%%%%%
\subsection{Edge Bound for Planar Graphs}
\label{subsec:planar-edge-bound}
%%%%%%%%%%%%%%%%%%%%%%%%%%%%%%%%%%%%%%%%%%%%%%%%%%%%%%%%%%%%

For a simple connected planar graph with

\[
|V|\geq3,
\]

every face has boundary length at least \(3\).

Counting edge-face incidences gives

\[
3|F|\leq2|E|.
\]

Using Euler's formula leads to

\[
\boxed{
|E|
\leq
3|V|-6.
}
\]

This provides a necessary condition for planarity.

%%%%%%%%%%%%%%%%%%%%%%%%%%%%%%%%%%%%%%%%%%%%%%%%%%%%%%%%%%%%
\subsection{Planar Bipartite Edge Bound}
\label{subsec:planar-bipartite-bound}
%%%%%%%%%%%%%%%%%%%%%%%%%%%%%%%%%%%%%%%%%%%%%%%%%%%%%%%%%%%%

A bipartite graph contains no triangle.

Thus every face of a simple planar bipartite graph has boundary length
at least \(4\).

Therefore,

\[
4|F|\leq2|E|.
\]

Combining with Euler's formula yields

\[
\boxed{
|E|
\leq
2|V|-4.
}
\]

%%%%%%%%%%%%%%%%%%%%%%%%%%%%%%%%%%%%%%%%%%%%%%%%%%%%%%%%%%%%
\subsection{Nonplanarity of \(K_5\)}
\label{subsec:k5-nonplanar}
%%%%%%%%%%%%%%%%%%%%%%%%%%%%%%%%%%%%%%%%%%%%%%%%%%%%%%%%%%%%

The complete graph \(K_5\) has

\[
|V|=5
\]

and

\[
|E|=10.
\]

If \(K_5\) were planar, the planar edge bound would require

\[
10\leq3(5)-6=9,
\]

which is false.

Therefore,

\[
\boxed{
K_5\text{ is nonplanar}.
}
\]

%%%%%%%%%%%%%%%%%%%%%%%%%%%%%%%%%%%%%%%%%%%%%%%%%%%%%%%%%%%%
\subsection{Nonplanarity of \(K_{3,3}\)}
\label{subsec:k33-nonplanar}
%%%%%%%%%%%%%%%%%%%%%%%%%%%%%%%%%%%%%%%%%%%%%%%%%%%%%%%%%%%%

The graph \(K_{3,3}\) is bipartite with

\[
|V|=6
\]

and

\[
|E|=9.
\]

If it were planar, the bipartite edge bound would require

\[
9\leq2(6)-4=8,
\]

which is impossible.

Therefore,

\[
\boxed{
K_{3,3}\text{ is nonplanar}.
}
\]

%%%%%%%%%%%%%%%%%%%%%%%%%%%%%%%%%%%%%%%%%%%%%%%%%%%%%%%%%%%%
\subsection{Kuratowski's Theorem}
\label{subsec:kuratowskis-theorem}
%%%%%%%%%%%%%%%%%%%%%%%%%%%%%%%%%%%%%%%%%%%%%%%%%%%%%%%%%%%%

\begin{theorem}[Kuratowski's Theorem]
A finite graph is planar if and only if it contains no subdivision of

\[
\boxed{
K_5
}
\]

or

\[
\boxed{
K_{3,3}.
}
\]
\end{theorem}

Thus two graphs characterize all obstructions to planarity through
subdivision.

%%%%%%%%%%%%%%%%%%%%%%%%%%%%%%%%%%%%%%%%%%%%%%%%%%%%%%%%%%%%
\subsection{Graph Traversal}
\label{subsec:graph-traversal}
%%%%%%%%%%%%%%%%%%%%%%%%%%%%%%%%%%%%%%%%%%%%%%%%%%%%%%%%%%%%

Graph traversal systematically visits vertices and edges.

Two fundamental algorithms are:

\[
\boxed{
\text{Breadth-First Search}
}
\]

and

\[
\boxed{
\text{Depth-First Search}.
}
\]

They form the foundation of many graph algorithms.

%%%%%%%%%%%%%%%%%%%%%%%%%%%%%%%%%%%%%%%%%%%%%%%%%%%%%%%%%%%%
\subsection{Breadth-First Search}
\label{subsec:breadth-first-search}
%%%%%%%%%%%%%%%%%%%%%%%%%%%%%%%%%%%%%%%%%%%%%%%%%%%%%%%%%%%%

Breadth-First Search, abbreviated BFS, explores vertices by increasing
distance from a starting vertex.

It uses a queue.

BFS can determine:

\begin{itemize}
    \item connected components;
    \item shortest paths in unweighted graphs;
    \item vertex distances;
    \item bipartiteness.
\end{itemize}

%%%%%%%%%%%%%%%%%%%%%%%%%%%%%%%%%%%%%%%%%%%%%%%%%%%%%%%%%%%%
\subsection{BFS Layers}
\label{subsec:bfs-layers}
%%%%%%%%%%%%%%%%%%%%%%%%%%%%%%%%%%%%%%%%%%%%%%%%%%%%%%%%%%%%

Starting from \(s\), define

\[
L_0=\{s\},
\]

and

\[
\boxed{
L_k
=
\{v:d(s,v)=k\}.
}
\]

BFS discovers these layers sequentially.

Thus the BFS tree records shortest-path distances from \(s\).

%%%%%%%%%%%%%%%%%%%%%%%%%%%%%%%%%%%%%%%%%%%%%%%%%%%%%%%%%%%%
\subsection{Depth-First Search}
\label{subsec:depth-first-search}
%%%%%%%%%%%%%%%%%%%%%%%%%%%%%%%%%%%%%%%%%%%%%%%%%%%%%%%%%%%%

Depth-First Search, abbreviated DFS, follows one branch as deeply as
possible before backtracking.

It uses a stack or recursion.

DFS can help identify:

\begin{itemize}
    \item connected components;
    \item cycles;
    \item articulation points;
    \item bridges;
    \item topological orderings in directed acyclic graphs.
\end{itemize}

%%%%%%%%%%%%%%%%%%%%%%%%%%%%%%%%%%%%%%%%%%%%%%%%%%%%%%%%%%%%
\subsection{Directed Acyclic Graphs}
\label{subsec:directed-acyclic-graphs}
%%%%%%%%%%%%%%%%%%%%%%%%%%%%%%%%%%%%%%%%%%%%%%%%%%%%%%%%%%%%

A directed acyclic graph, abbreviated DAG, is a directed graph containing
no directed cycle.

DAGs naturally represent precedence structures.

Examples include:

\begin{itemize}
    \item prerequisite networks;
    \item task dependencies;
    \item build systems;
    \item scheduling constraints.
\end{itemize}

%%%%%%%%%%%%%%%%%%%%%%%%%%%%%%%%%%%%%%%%%%%%%%%%%%%%%%%%%%%%
\subsection{Topological Ordering}
\label{subsec:topological-ordering}
%%%%%%%%%%%%%%%%%%%%%%%%%%%%%%%%%%%%%%%%%%%%%%%%%%%%%%%%%%%%

A topological ordering of a DAG is a linear ordering of its vertices
such that every directed edge

\[
u\to v
\]

places \(u\) before \(v\).

\begin{theorem}
A finite directed graph has a topological ordering if and only if it is
acyclic.
\end{theorem}

%%%%%%%%%%%%%%%%%%%%%%%%%%%%%%%%%%%%%%%%%%%%%%%%%%%%%%%%%%%%
\subsection{Shortest Paths}
\label{subsec:shortest-paths}
%%%%%%%%%%%%%%%%%%%%%%%%%%%%%%%%%%%%%%%%%%%%%%%%%%%%%%%%%%%%

In an unweighted graph, BFS finds shortest paths.

In a weighted graph with nonnegative edge weights, Dijkstra's algorithm
finds shortest paths from one source.

The path length is

\[
\boxed{
\sum_{e\in P}w(e).
}
\]

Shortest-path problems are fundamental in transportation and network
optimization.

%%%%%%%%%%%%%%%%%%%%%%%%%%%%%%%%%%%%%%%%%%%%%%%%%%%%%%%%%%%%
\subsection{Minimum Spanning Trees}
\label{subsec:minimum-spanning-trees}
%%%%%%%%%%%%%%%%%%%%%%%%%%%%%%%%%%%%%%%%%%%%%%%%%%%%%%%%%%%%

Suppose a connected graph has weighted edges.

A minimum spanning tree is a spanning tree minimizing

\[
\boxed{
\sum_{e\in T}w(e).
}
\]

Two classical algorithms are:

\begin{itemize}
    \item Kruskal's algorithm;
    \item Prim's algorithm.
\end{itemize}

Minimum spanning trees connect graph theory with greedy optimization.

%%%%%%%%%%%%%%%%%%%%%%%%%%%%%%%%%%%%%%%%%%%%%%%%%%%%%%%%%%%%
\subsection{Kruskal's Algorithm}
\label{subsec:kruskals-algorithm}
%%%%%%%%%%%%%%%%%%%%%%%%%%%%%%%%%%%%%%%%%%%%%%%%%%%%%%%%%%%%

Kruskal's algorithm repeatedly selects the currently cheapest edge that
does not create a cycle.

The process stops after selecting

\[
|V|-1
\]

edges.

The result is a minimum spanning tree.

%%%%%%%%%%%%%%%%%%%%%%%%%%%%%%%%%%%%%%%%%%%%%%%%%%%%%%%%%%%%
\subsection{Prim's Algorithm}
\label{subsec:prims-algorithm}
%%%%%%%%%%%%%%%%%%%%%%%%%%%%%%%%%%%%%%%%%%%%%%%%%%%%%%%%%%%%

Prim's algorithm grows one connected tree.

Starting from one vertex, repeatedly add the cheapest edge joining the
current tree to a new vertex.

When all vertices are included, the resulting tree is minimum-weight.

%%%%%%%%%%%%%%%%%%%%%%%%%%%%%%%%%%%%%%%%%%%%%%%%%%%%%%%%%%%%
\subsection{Graph Laplacian}
\label{subsec:graph-laplacian}
%%%%%%%%%%%%%%%%%%%%%%%%%%%%%%%%%%%%%%%%%%%%%%%%%%%%%%%%%%%%

Let \(A\) be the adjacency matrix and \(D\) the diagonal degree matrix:

\[
D_{ii}=\deg(v_i).
\]

The graph Laplacian is

\[
\boxed{
L=D-A.
}
\]

The Laplacian connects graph structure with linear algebra,
probability, geometry, and spectral theory.

%%%%%%%%%%%%%%%%%%%%%%%%%%%%%%%%%%%%%%%%%%%%%%%%%%%%%%%%%%%%
\subsection{Laplacian and Connectivity}
\label{subsec:laplacian-connectivity}
%%%%%%%%%%%%%%%%%%%%%%%%%%%%%%%%%%%%%%%%%%%%%%%%%%%%%%%%%%%%

The vector

\[
\mathbf{1}
=
(1,\ldots,1)^T
\]

satisfies

\[
L\mathbf{1}=0.
\]

The multiplicity of the eigenvalue \(0\) equals the number of connected
components.

Thus,

\[
\boxed{
G\text{ connected}
\iff
0\text{ is a simple Laplacian eigenvalue}.
}
\]

%%%%%%%%%%%%%%%%%%%%%%%%%%%%%%%%%%%%%%%%%%%%%%%%%%%%%%%%%%%%
\subsection{Matrix--Tree Theorem}
\label{subsec:matrix-tree-theorem}
%%%%%%%%%%%%%%%%%%%%%%%%%%%%%%%%%%%%%%%%%%%%%%%%%%%%%%%%%%%%

\begin{theorem}[Matrix--Tree Theorem]
The number of spanning trees of a finite graph \(G\) equals any cofactor
of its Laplacian matrix.
\end{theorem}

That is, delete one row and the corresponding column from \(L\).

The determinant of the resulting matrix is

\[
\boxed{
\tau(G),
}
\]

the number of spanning trees.

This provides a striking connection between determinants and
combinatorial enumeration.

%%%%%%%%%%%%%%%%%%%%%%%%%%%%%%%%%%%%%%%%%%%%%%%%%%%%%%%%%%%%
\subsection{Network Flow Preview}
\label{subsec:network-flow-preview}
%%%%%%%%%%%%%%%%%%%%%%%%%%%%%%%%%%%%%%%%%%%%%%%%%%%%%%%%%%%%

A flow network is a directed graph in which edges have capacities.

One seeks to send as much flow as possible from a source \(s\) to a
sink \(t\).

The central result is the max-flow min-cut theorem.

%%%%%%%%%%%%%%%%%%%%%%%%%%%%%%%%%%%%%%%%%%%%%%%%%%%%%%%%%%%%
\subsection{Max-Flow Min-Cut Theorem}
\label{subsec:max-flow-min-cut-preview}
%%%%%%%%%%%%%%%%%%%%%%%%%%%%%%%%%%%%%%%%%%%%%%%%%%%%%%%%%%%%

\begin{theorem}[Max-Flow Min-Cut Theorem]
In a finite capacitated network,

\[
\boxed{
\text{maximum }s\text{-}t\text{ flow}
=
\text{minimum }s\text{-}t\text{ cut capacity}.
}
\]
\end{theorem}

This is another fundamental min-max principle.

Network flow is developed more fully in optimization and operations
research.

%%%%%%%%%%%%%%%%%%%%%%%%%%%%%%%%%%%%%%%%%%%%%%%%%%%%%%%%%%%%
\subsection{Graph Products Preview}
\label{subsec:graph-products-preview}
%%%%%%%%%%%%%%%%%%%%%%%%%%%%%%%%%%%%%%%%%%%%%%%%%%%%%%%%%%%%

Graphs may be combined through several product constructions.

Examples include:

\begin{itemize}
    \item Cartesian product;
    \item direct product;
    \item strong product;
    \item lexicographic product.
\end{itemize}

For example, the \(n\)-dimensional hypercube satisfies

\[
\boxed{
Q_n
=
K_2\square K_2\square\cdots\square K_2.
}
\]

The symbol

\[
\square
\]

denotes the Cartesian graph product.

%%%%%%%%%%%%%%%%%%%%%%%%%%%%%%%%%%%%%%%%%%%%%%%%%%%%%%%%%%%%
\subsection{Hypercube Graphs}
\label{subsec:hypercube-graphs}
%%%%%%%%%%%%%%%%%%%%%%%%%%%%%%%%%%%%%%%%%%%%%%%%%%%%%%%%%%%%

The hypercube graph \(Q_n\) has vertex set

\[
\{0,1\}^n.
\]

Two binary strings are adjacent when they differ in exactly one
coordinate.

Thus adjacency corresponds to Hamming distance \(1\).

The graph has

\[
\boxed{
2^n
}
\]

vertices.

Every vertex has degree

\[
n.
\]

Hence,

\[
\boxed{
|E(Q_n)|
=
n2^{n-1}.
}
\]

%%%%%%%%%%%%%%%%%%%%%%%%%%%%%%%%%%%%%%%%%%%%%%%%%%%%%%%%%%%%
\subsection{Graph Homomorphisms}
\label{subsec:graph-homomorphisms}
%%%%%%%%%%%%%%%%%%%%%%%%%%%%%%%%%%%%%%%%%%%%%%%%%%%%%%%%%%%%

A graph homomorphism from \(G\) to \(H\) is a map

\[
f:V(G)\to V(H)
\]

such that adjacency is preserved:

\[
\boxed{
uv\in E(G)
\Longrightarrow
f(u)f(v)\in E(H).
}
\]

Proper \(k\)-colorings are exactly graph homomorphisms

\[
G\to K_k.
\]

Thus coloring may be viewed as a structure-preserving map problem.

%%%%%%%%%%%%%%%%%%%%%%%%%%%%%%%%%%%%%%%%%%%%%%%%%%%%%%%%%%%%
\subsection{Applications}
\label{subsec:graph-theory-applications}
%%%%%%%%%%%%%%%%%%%%%%%%%%%%%%%%%%%%%%%%%%%%%%%%%%%%%%%%%%%%

\begin{applicationbox}

\textbf{Transportation.}
Roads, airline routes, shipping lanes, and transit systems are modeled
as graphs.

\medskip

\textbf{Computer Networks.}
Routers, devices, and communication links form graph structures.

\medskip

\textbf{Social Networks.}
People or organizations form vertices while relationships form edges.

\medskip

\textbf{Biology.}
Protein interactions, food webs, neural networks, and phylogenetic
relationships use graph models.

\medskip

\textbf{Chemistry.}
Molecules can be represented as graphs with atoms as vertices and bonds
as edges.

\medskip

\textbf{Scheduling.}
Precedence constraints form DAGs, while incompatibility constraints can
form coloring graphs.

\medskip

\textbf{Optimization.}
Shortest paths, spanning trees, matchings, flows, cuts, and routing are
graph optimization problems.

\medskip

\textbf{Data Science.}
Graph models describe relational, networked, and non-Euclidean data.

\medskip

\textbf{Probability.}
Random graph theory studies typical network structure and threshold
phenomena.

\medskip

\textbf{Algebra.}
Spectral graph theory uses matrices, eigenvalues, and algebraic
invariants to study graphs.

\end{applicationbox}

%%%%%%%%%%%%%%%%%%%%%%%%%%%%%%%%%%%%%%%%%%%%%%%%%%%%%%%%%%%%
\subsection{Common Errors}
\label{subsec:graph-theory-common-errors}
%%%%%%%%%%%%%%%%%%%%%%%%%%%%%%%%%%%%%%%%%%%%%%%%%%%%%%%%%%%%

\subsubsection{Confusing a Graph with Its Drawing}

Different drawings may represent the same graph.

Edge crossings in a drawing do not create vertices unless explicitly
declared.

%%%%%%%%%%%%%%%%%%%%%%%%%%%%%%%%%%%%%%%%%%%%%%%%%%%%%%%%%%%%
\subsubsection{Confusing Walks, Trails, and Paths}

A walk may repeat vertices and edges.

A trail cannot repeat edges.

A path cannot repeat vertices.

%%%%%%%%%%%%%%%%%%%%%%%%%%%%%%%%%%%%%%%%%%%%%%%%%%%%%%%%%%%%
\subsubsection{Confusing Euler and Hamiltonian Problems}

Euler problems concern using every edge.

Hamiltonian problems concern visiting every vertex.

These are fundamentally different conditions.

%%%%%%%%%%%%%%%%%%%%%%%%%%%%%%%%%%%%%%%%%%%%%%%%%%%%%%%%%%%%
\subsubsection{Using Degree Sum Without Dividing by Two}

Because each edge contributes to two vertex degrees,

\[
|E|
=
\frac12
\sum_{v\in V}\deg(v).
\]

%%%%%%%%%%%%%%%%%%%%%%%%%%%%%%%%%%%%%%%%%%%%%%%%%%%%%%%%%%%%
\subsubsection{Assuming Every Degree Sequence Is Graphical}

A list of nonnegative integers may fail to represent any simple graph.

For example, degree values may violate structural constraints even when
their total sum is even.

%%%%%%%%%%%%%%%%%%%%%%%%%%%%%%%%%%%%%%%%%%%%%%%%%%%%%%%%%%%%
\subsubsection{Confusing Connected and Complete}

Connected means every pair is joined by some path.

Complete means every pair is joined directly by an edge.

Every complete graph is connected, but most connected graphs are not
complete.

%%%%%%%%%%%%%%%%%%%%%%%%%%%%%%%%%%%%%%%%%%%%%%%%%%%%%%%%%%%%
\subsubsection{Assuming Every Connected Graph Is a Tree}

A tree must be connected and acyclic.

Connected graphs may contain many cycles.

%%%%%%%%%%%%%%%%%%%%%%%%%%%%%%%%%%%%%%%%%%%%%%%%%%%%%%%%%%%%
\subsubsection{Using Euler's Planar Formula for a Nonplanar Drawing}

Euler's formula applies to a planar embedding.

A drawing with crossings is not automatically a planar embedding.

%%%%%%%%%%%%%%%%%%%%%%%%%%%%%%%%%%%%%%%%%%%%%%%%%%%%%%%%%%%%
\subsubsection{Assuming the Planar Edge Bound Proves Planarity}

The inequality

\[
|E|\leq3|V|-6
\]

is necessary for simple planar graphs, but it is not sufficient.

%%%%%%%%%%%%%%%%%%%%%%%%%%%%%%%%%%%%%%%%%%%%%%%%%%%%%%%%%%%%
\subsubsection{Confusing Matching and Vertex Cover}

A matching is a set of pairwise disjoint edges.

A vertex cover is a set of vertices touching every edge.

Their relationship is especially strong in bipartite graphs through
K\H{o}nig's theorem.

%%%%%%%%%%%%%%%%%%%%%%%%%%%%%%%%%%%%%%%%%%%%%%%%%%%%%%%%%%%%
\subsection{Exercises}
\label{subsec:graph-theory-exercises}
%%%%%%%%%%%%%%%%%%%%%%%%%%%%%%%%%%%%%%%%%%%%%%%%%%%%%%%%%%%%

\begin{exercise}[1]
Define a simple graph.
\end{exercise}

\begin{exercise}[1]
What is the maximum number of edges in a simple graph on \(8\)
vertices?
\end{exercise}

\begin{exercise}[1]
Find the number of edges in

\[
K_7.
\]
\end{exercise}

\begin{exercise}[1]
Find the number of edges in

\[
K_{4,6}.
\]
\end{exercise}

\begin{exercise}[1]
A graph has degree sequence

\[
4,4,3,3,2,2.
\]

Find the number of edges.
\end{exercise}

\begin{exercise}[1]
Explain why a graph cannot have exactly three odd-degree vertices.
\end{exercise}

\begin{exercise}[1]
Find the degree of every vertex in

\[
C_8.
\]
\end{exercise}

\begin{exercise}[1]
Find the number of edges in the \(5\)-dimensional hypercube \(Q_5\).
\end{exercise}

\begin{exercise}[1]
Define a connected graph.
\end{exercise}

\begin{exercise}[1]
Define a tree.
\end{exercise}

\begin{exercise}[1]
How many edges does a tree on \(20\) vertices have?
\end{exercise}

\begin{exercise}[1]
How many edges does a forest with \(20\) vertices and \(4\) components
have?
\end{exercise}

\begin{exercise}[1]
State the Euler circuit criterion.
\end{exercise}

\begin{exercise}[1]
State the Euler trail criterion.
\end{exercise}

\begin{exercise}[1]
Define the chromatic number.
\end{exercise}

\begin{exercise}[1]
Find

\[
\chi(K_6).
\]
\end{exercise}

\begin{exercise}[1]
Find

\[
\chi(C_8).
\]
\end{exercise}

\begin{exercise}[1]
Find

\[
\chi(C_7).
\]
\end{exercise}

\begin{exercise}[2]
Prove the handshaking lemma by double counting incidences.
\end{exercise}

\begin{exercise}[2]
Prove that every graph has an even number of odd-degree vertices.
\end{exercise}

\begin{exercise}[2]
Determine whether

\[
(3,3,2,2,2,2)
\]

is graphical.
\end{exercise}

\begin{exercise}[2]
Use Havel--Hakimi to test

\[
(4,4,3,2,2,1).
\]
\end{exercise}

\begin{exercise}[2]
Prove that every tree with at least two vertices has at least two
leaves.
\end{exercise}

\begin{exercise}[2]
Prove that every tree on \(n\) vertices has \(n-1\) edges.
\end{exercise}

\begin{exercise}[2]
Prove that every connected graph contains a spanning tree.
\end{exercise}

\begin{exercise}[2]
Find the number of labeled trees on \(5\) vertices using Cayley's
formula.
\end{exercise}

\begin{exercise}[2]
List all Prüfer sequences for labeled trees on \(3\) vertices.
\end{exercise}

\begin{exercise}[2]
Determine whether a connected graph with degrees

\[
4,4,2,2,2,2
\]

has an Euler circuit.
\end{exercise}

\begin{exercise}[2]
Determine whether a connected graph with degrees

\[
3,3,2,2,2,2
\]

has an Euler trail.
\end{exercise}

\begin{exercise}[2]
Show that every bipartite graph contains no odd cycle.
\end{exercise}

\begin{exercise}[2]
Show that every tree is bipartite.
\end{exercise}

\begin{exercise}[2]
Find the adjacency matrix of the path graph \(P_4\).
\end{exercise}

\begin{exercise}[2]
Square the adjacency matrix of \(P_4\) and interpret one nonzero entry.
\end{exercise}

\begin{exercise}[2]
Use Euler's planar formula to determine the number of faces in a
connected planar graph with \(10\) vertices and \(15\) edges.
\end{exercise}

\begin{exercise}[2]
Use the planar edge bound to prove that \(K_5\) is nonplanar.
\end{exercise}

\begin{exercise}[2]
Use the bipartite planar edge bound to prove that \(K_{3,3}\) is
nonplanar.
\end{exercise}

\begin{exercise}[2]
State Hall's Marriage Theorem.
\end{exercise}

\begin{exercise}[2]
Give a bipartite graph satisfying Hall's condition and construct a
matching saturating the left part.
\end{exercise}

\begin{exercise}[3]
Prove that a connected graph with \(n\) vertices and \(n-1\) edges is a
tree.
\end{exercise}

\begin{exercise}[3]
Prove that an acyclic graph with \(n\) vertices and \(n-1\) edges is
connected.
\end{exercise}

\begin{exercise}[3]
Prove that an edge is a bridge if and only if it belongs to no cycle.
\end{exercise}

\begin{exercise}[3]
Prove that a graph is bipartite if and only if it contains no odd
cycle.
\end{exercise}

\begin{exercise}[3]
Derive the formula

\[
|E(Q_n)|=n2^{n-1}.
\]
\end{exercise}

\begin{exercise}[3]
Prove that

\[
P_T(k)=k(k-1)^{n-1}
\]

for every tree \(T\) on \(n\) vertices.
\end{exercise}

\begin{exercise}[3]
Verify deletion--contraction for the chromatic polynomial of \(C_3\).
\end{exercise}

\begin{exercise}[3]
Prove that

\[
\chi(G)\leq\Delta(G)+1
\]

using greedy coloring.
\end{exercise}

\begin{exercise}[3]
Prove that

\[
\alpha(G)+\tau(G)=|V|.
\]
\end{exercise}

\begin{exercise}[3]
Apply Hall's theorem to determine whether a perfect assignment exists
in a bipartite graph of your own construction.
\end{exercise}

\begin{exercise}[3]
Prove the planar edge bound

\[
|E|\leq3|V|-6.
\]
\end{exercise}

\begin{exercise}[3]
Prove the bipartite planar edge bound

\[
|E|\leq2|V|-4.
\]
\end{exercise}

\begin{exercise}[3]
Show that the multiplicity of the Laplacian eigenvalue \(0\) equals the
number of connected components for a small example.
\end{exercise}

\begin{exercise}[3]
Use the Matrix--Tree Theorem to compute the number of spanning trees of
\(K_4\).
\end{exercise}

\begin{exercise}[3]
Use BFS to determine shortest-path distances from a chosen vertex in a
graph of your own construction.
\end{exercise}

\begin{exercise}[3]
Use DFS to identify the connected components of a disconnected graph.
\end{exercise}

\begin{exercise}[4]
Study the Erd\H{o}s--Gallai theorem for graphical degree sequences.
\end{exercise}

\begin{exercise}[4]
Investigate Menger's theorem and its relationship to vertex and edge
connectivity.
\end{exercise}

\begin{exercise}[4]
Study the Matrix--Tree Theorem and prove it using Laplacian
determinants.
\end{exercise}

\begin{exercise}[4]
Investigate spectral graph theory and the relationship between adjacency
eigenvalues and graph structure.
\end{exercise}

\begin{exercise}[4]
Study Brooks' theorem as a refinement of the greedy coloring bound.
\end{exercise}

\begin{exercise}[4]
Investigate the Four Color Theorem.
\end{exercise}

\begin{exercise}[4]
Study Tutte's theorem characterizing graphs with perfect matchings.
\end{exercise}

\begin{exercise}[4]
Investigate K\H{o}nig's theorem and prove it using augmenting paths.
\end{exercise}

\begin{exercise}[4]
Study the max-flow min-cut theorem and its relationship to bipartite
matching.
\end{exercise}

\begin{exercise}[4]
Investigate Kuratowski's theorem and Wagner's theorem on planarity.
\end{exercise}

%%%%%%%%%%%%%%%%%%%%%%%%%%%%%%%%%%%%%%%%%%%%%%%%%%%%%%%%%%%%
\subsection{Conceptual Summary}
\label{subsec:graph-theory-summary}
%%%%%%%%%%%%%%%%%%%%%%%%%%%%%%%%%%%%%%%%%%%%%%%%%%%%%%%%%%%%

A graph is a pair

\[
\boxed{
G=(V,E)
}
\]

representing vertices and pairwise relationships.

The degree of a vertex satisfies the handshaking identity

\[
\boxed{
\sum_{v\in V}\deg(v)=2|E|.
}
\]

Important graph classes include

\[
\boxed{
K_n,
\quad
K_{m,n},
\quad
C_n,
\quad
P_n,
\quad
Q_n.
}
\]

A graph is connected when every pair of vertices is joined by a path.

A tree is a connected acyclic graph and satisfies

\[
\boxed{
|E|=|V|-1.
}
\]

A forest with \(c\) components satisfies

\[
\boxed{
|E|=|V|-c.
}
\]

Euler circuits are characterized by even vertex degrees.

Hamiltonian cycles visit every vertex exactly once.

The chromatic number

\[
\boxed{
\chi(G)
}
\]

measures the minimum number of colors needed for a proper coloring.

Matchings pair vertices without overlap.

Hall's theorem characterizes when a bipartite matching can saturate one
part.

For bipartite graphs,

\[
\boxed{
\nu(G)=\tau(G)
}
\]

by K\H{o}nig's theorem.

Planar graphs satisfy Euler's relation

\[
\boxed{
|V|-|E|+|F|=2.
}
\]

Simple planar graphs satisfy

\[
\boxed{
|E|\leq3|V|-6.
}
\]

Adjacency matrices encode edges, and

\[
\boxed{
(A^k)_{ij}
}
\]

counts length-\(k\) walks.

The Laplacian

\[
\boxed{
L=D-A
}
\]

encodes connectivity and spanning-tree information.

Graph theory therefore provides a unified mathematical language for

\[
\boxed{
\text{networks}
+
\text{paths}
+
\text{connectivity}
+
\text{coloring}
+
\text{matching}
+
\text{optimization}.
}
\]

%%%%%%%%%%%%%%%%%%%%%%%%%%%%%%%%%%%%%%%%%%%%%%%%%%%%%%%%%%%%
\subsection{Section Connections}
\label{subsec:graph-theory-connections}
%%%%%%%%%%%%%%%%%%%%%%%%%%%%%%%%%%%%%%%%%%%%%%%%%%%%%%%%%%%%

\chapterconnection{
Graph Theory builds on the discrete structures introduced in
Section~\ref{sec:discrete-structures} and uses the counting,
enumerative, extremal, algebraic, and probabilistic methods developed
throughout Chapter 6. Extremal Combinatorics studies edge bounds under
forbidden-subgraph conditions. Algebraic Combinatorics contributes
chromatic polynomials, spectra, group actions, and the Tutte polynomial.
Probabilistic Combinatorics studies random graphs and typical graph
behavior. The next section, Ramsey Theory, asks when large graphs or
edge colorings must contain unavoidable configurations. Design Theory
uses highly regular incidence graphs, while Matroid Theory generalizes
trees, independence, rank, circuits, and deletion--contraction.
Optimization and Operations Research develops shortest paths, flows,
matchings, network optimization, and spanning-tree methods in greater
algorithmic depth.
}

%%%%%%%%%%%%%%%%%%%%%%%%%%%%%%%%%%%%%%%%%%%%%%%%%%%%%%%%%%%%
% SECTION SOURCE TRACKING
%%%%%%%%%%%%%%%%%%%%%%%%%%%%%%%%%%%%%%%%%%%%%%%%%%%%%%%%%%%%

%------------------------------------------------------------
% 6.9 Graph Theory
%------------------------------------------------------------
%
% Source basis:
%   - Standard introductory graph theory.
%   - Standard algebraic and algorithmic graph theory.
%   - Standard network and connectivity theory.
%
% Used for:
%   - graphs and graph types;
%   - vertices and edges;
%   - simple graphs;
%   - multigraphs and pseudographs;
%   - directed and weighted graphs;
%   - degree and neighborhoods;
%   - handshaking lemma;
%   - regular graphs;
%   - complete and bipartite graphs;
%   - multipartite graphs;
%   - subgraphs and induced subgraphs;
%   - complements;
%   - graph isomorphism;
%   - invariants and degree sequences;
%   - Havel--Hakimi method;
%   - walks, trails, paths, circuits, and cycles;
%   - distance, diameter, radius, and eccentricity;
%   - connectivity and components;
%   - cut vertices and bridges;
%   - vertex and edge connectivity;
%   - adjacency and incidence matrices;
%   - adjacency-matrix walk counting;
%   - trees and forests;
%   - spanning trees;
%   - Cayley's formula;
%   - Prüfer codes;
%   - rooted trees;
%   - Euler trails and circuits;
%   - Hamiltonian paths and cycles;
%   - Dirac's theorem;
%   - Ore's theorem;
%   - vertex and edge coloring;
%   - chromatic number;
%   - greedy coloring;
%   - chromatic polynomial;
%   - deletion--contraction;
%   - matchings;
%   - Hall's theorem;
%   - vertex covers and independent sets;
%   - König's theorem;
%   - planar graphs;
%   - Euler's planar formula;
%   - planar edge bounds;
%   - K5 and K3,3 nonplanarity;
%   - Kuratowski's theorem;
%   - BFS and DFS;
%   - DAGs and topological ordering;
%   - shortest paths;
%   - minimum spanning trees;
%   - Kruskal and Prim algorithms;
%   - graph Laplacian;
%   - Matrix--Tree Theorem;
%   - network flow preview;
%   - graph products;
%   - hypercube graphs;
%   - graph homomorphisms;
%   - applications and exercises.
%
% Notes:
%   - Random graph theory is introduced in Section 6.8.
%   - Ramsey Theory follows this section and studies unavoidable
%     monochromatic and structural configurations.
%   - Design Theory and Matroid Theory reuse incidence, matching,
%     independence, rank, and graph-structural ideas.
%   - Network optimization algorithms are developed further in
%     Chapter 10.
%
%%%%%%%%%%%%%%%%%%%%%%%%%%%%%%%%%%%%%%%%%%%%%%%%%%%%%%%%%%%%